\documentclass[a4paper,11pt]{article}
\usepackage[margin=27mm]{geometry}
\usepackage{fontspec,amsmath,amssymb,amsthm,unicode-math}
\setmainfont{STIXGeneral}\setmathfont{STIX Two Math}
\newtheorem{theorem}{Theorem}\newtheorem{lemma}[theorem]{Lemma}
\newtheorem{corollary}[theorem]{Corollary}\newtheorem{proposition}[theorem]{Proposition}
\newcommand{\Aut}{\operatorname{Aut}}\newcommand{\Stab}{\operatorname{Stab}}\newcommand{\Pos}{\operatorname{Pos}}
\newcommand{\Fl}{\operatorname{Fl}}\newcommand{\gen}[1]{\langle#1\rangle}
\title{Equivalence relations, transporters and flags}
\author{Human Compression project}\date{Working proof edition; independent review pending}
\begin{document}\maketitle
\begin{abstract}
A family of equivalence relations has a common refinement. If a group
preserves the relations, each part of the refinement meets a group orbit
in either an empty set or a torsor for an intersection of stabilizers.
If maps preserve the relations and jointly distinguish points, faithfulness
passes from their targets to their source. We give both statements explicitly
and apply them to the faces of colored graphs. This supplies the symmetry,
covering and mixing criteria of Cunningham, Moch\'an, Montero, Poto\v cnik and
Toledo within one elementary framework. We also give the precise connectivity
criterion for the natural reconstruction of a thin poset from its flag graph.
The transporter and group-action arguments are standard; their organization
here is not a claim of historical priority.
\end{abstract}

\section{The common refinement}
Let $X$ be a set with equivalence relations $E_i$, indexed by a set $I$.
Write $x\sim y$ when $xE_i y$ for every $i$. Its class is
\[
 C(x)=\bigcap_{i\in I}[x]_{E_i}.
\]
We call the family separating if every $C(x)$ is a singleton. Empty sets
of indices are allowed, although the graph applications have positive rank.

\begin{theorem}\label{transport}
Suppose a group $G$ acts on the right on $X$ and preserves every $E_i$.
Put $H_i(x)=\{g:xgE_i x\}$ and $H(x)=\bigcap_i H_i(x)$.
For $x,y\in X$, put $T_i(x,y)=\{g:ygE_i x\}$.
Then $H_i(x)$ is the stabilizer of the class $[x]_{E_i}$, and
\[
 \{g:yg\sim x\}=\bigcap_i T_i(x,y).
 \tag{1}
\]
Each nonempty $T_i(x,y)$ equals $\tau_i H_i(x)$ for any
$\tau_i\in T_i(x,y)$. If the intersection in (1) is nonempty, it equals
$g_0H(x)$ for any element $g_0$ of that intersection. Consequently
\[
 C(x)\cap yG=yg_0H(x).
 \tag{2}
\]
If the action is free, every nonempty set in (2) is an $H(x)$-torsor.
In particular, if $\Omega$ contains one representative of each $G$-orbit
meeting $C(x)$, chosen within that class, then
\[
 C(x)=\coprod_{y\in\Omega}yH(x).
 \tag{3}
\]
\end{theorem}
\begin{proof}
Preservation of an equivalence relation means that $[xg]_{E_i}$ is the
image of $[x]_{E_i}$ under $g$. Thus $xgE_i x$ is exactly the condition
that $g$ stabilize this class, proving the subgroup assertion.
The first equality is the definition of the common refinement.
If $y\tau_iE_i x$, apply $\tau_i^{-1}g$ to this relation. It gives
$ygE_i x(\tau_i^{-1}g)$, so $ygE_i x$ holds precisely when
$\tau_i^{-1}g\in H_i(x)$. This proves the transporter formula.
If $g_0$ lies in every transporter, rewrite each one as $g_0H_i(x)$
and intersect. Evaluation at $y$ gives (2). In a free action its
restriction to $g_0H(x)$ is injective, as well as surjective onto (2).
Applying this to the distinct orbits gives the disjoint union (3).
No finiteness assumption occurs in the argument.
\end{proof}

For a free action, the family is separating exactly when $H(x)=1$
at one representative of every orbit and all the intersections (1)
between distinct orbits are empty. Without freeness the exact condition
is $H(x)=G_x$, together with those empty intersections, where $G_x$
is the point stabilizer. This distinction prevents an action kernel
from being mistaken for a failure of separation.

\begin{theorem}\label{joint}
Let $f_a:X\to Y_a$ preserve each indexed equivalence relation, and
suppose their joint map $X\to\prod_aY_a$ is injective. If every target
family is separating, then so is the family on $X$.
More generally, if one target family is separating and $f$ is injective
on each class of one of the relations $E_j$, then the family on $X$
is separating.
\end{theorem}
\begin{proof}
If $x\sim x'$, then $f_a(x)\sim f_a(x')$ for every $a$. Separation
in the targets makes these images equal; joint injectivity gives $x=x'$.
For the second assertion, separation gives $f(x)=f(x')$, while
$xE_jx'$ puts the two points in a class on which $f$ is injective.
\end{proof}

These statements are the same fiber method as the intersection theorem
for simultaneous group actions, applied to the map recording all faces
rather than to two subgroup orbits. Formula (1) retains the actual
cosets; replacing their intersection by an intersection of abstractly
isomorphic groups loses the relative translations and is generally false.

\section{Faces of a colored graph}
Let $n\geq2$. An $n$-premaniplex is a connected set of flags with
involutions $r_0,\ldots,r_{n-1}$ satisfying $r_ir_j=r_jr_i$ when
$|i-j|>1$. A maniplex additionally has no fixed points of any $r_i$
or $r_ir_j$ with $i\ne j$. We use right actions. For
$J\subseteq\{0,\ldots,n-1\}$, write $W_J=\gen{r_j:j\in J}$,
interpreted as a subgroup of the universal string group acting on the flags.
The $i$-face at $x$ is
\[
 F_i(x)=xW_{I\setminus\{i\}}.
\]
Thus $E_i$ means membership in the same $i$-face, and $C(x)$ consists
of flags having every face in common with $x$.

We recall why this agrees with the face-poset definition. Faces of
different ranks are incident when they meet. Their order is transitive,
and every chain is realized by a flag. To check both assertions, split
a path avoiding color $k$ into its colors below $k$ followed by its
colors above $k$. These two sets commute. If a flag already lies in
prescribed faces of ranks at least $k$, move toward an incident face
of rank $j<k$ along a path in the prescribed $k$-face. After this
reordering, its lower-color part preserves every prescribed higher
face, and its upper-color part preserves the new $j$-face. The flag
between the two parts lies in all these faces. Three ranks give
transitivity; descending induction gives realization of a full chain.
This is the argument used in \emph{Intersections and faces in graph
covers}; it does not require finite flag sets.

The natural map from flags to maximal chains is therefore surjective,
with fibers $C(x)$. A maniplex is faithful precisely when its face
relations are separating. An automorphism commutes with every color
involution, so preserves these relations. Its action is free: fixing
one flag fixes every flag reached by a color word, hence the whole
connected graph. Theorem~\ref{transport} therefore applies to every
subgroup of its automorphism group.

For completeness, the tests between different automorphism orbits can
be made on their quotient graph. If their quotient flags are not
equivalent for some $E_i$, their lifts cannot be chain-equivalent.
Otherwise choose an $i$-avoiding path in the quotient and lift it
from $x$. It ends in the orbit of $y$, at some $y\tau_i$, giving
the transporter in (1). Path lifting requires only the unique edge
of each color, and does not presume that all path lifts have the
same endpoint.

\section{Deleting one face}
Write $C_{\widehat i}(x)=\bigcap_{j\ne i}F_j(x)$. The following
elementary form of the diamond condition is useful because it survives
projection.
\begin{lemma}\label{thin}
A faithful maniplex is thin if and only if, for every $x$ and $i$,
\[
 C_{\widehat i}(x)=\{x,xr_i\}.
 \tag{4}
\]
Here thinness includes intervals at either improper endpoint.
\end{lemma}
\begin{proof}
The two displayed flags share every face except the $i$-face, and
those two $i$-faces are distinct by faithfulness. If the face poset
is thin, there are exactly two completions of the chain missing rank
$i$. Realization and faithfulness identify their flags with $x$ and
$xr_i$, proving (4).

Conversely fix incident faces of ranks $i-1$ and $i+1$, using the
improper endpoint when necessary. Extend them to a maximal chain
and realize that chain by $x$. Any intermediate $i$-face is comparable
with every other member of this chain, by transitivity of the order.
It therefore completes the chain missing rank $i$. Its realization
lies in $C_{\widehat i}(x)$ and hence is one of the two flags in (4).
There are exactly two intermediate faces, as required.
\end{proof}

\begin{theorem}\label{mix}
Let $(X_a,x_a)$ be a nonempty family of rooted faithful maniplexes
of the same rank. The orbit of $(x_a)_a$ under simultaneous color
adjacency is faithful. If every $X_a$ is also thin and orientable,
then this orbit is thin and orientable.
\end{theorem}
\begin{proof}
The coordinate projections are covers and preserve every face relation.
They jointly distinguish tuples, so Theorem~\ref{joint} proves
faithfulness. Any coordinate shows that the simultaneous orbit has
no loops or multiple edges, and its distant colors commute. Thus it
is a maniplex, for finite or infinite index sets alike.

Suppose the factors are thin and orientable. Choose the two-coloring
of each factor that assigns color zero to its root. A color word
changes every coordinate's color by the parity of its length.
Every tuple in the orbit consequently has the same color in all
coordinates. For $q=(z_a)_a$, project a flag of
$C_{\widehat i}(q)$ to a coordinate. Equation (4) puts that coordinate
in $\{z_a,z_ar_i\}$. The two alternatives have different colors.
Equality of colors across coordinates therefore forces the same
alternative everywhere. Hence
$C_{\widehat i}(q)=\{q,qr_i\}$, and Lemma~\ref{thin} proves thinness.
Any coordinate coloring gives orientability of the simultaneous orbit.
\end{proof}

\begin{corollary}\label{regularcover}
The smallest regular cover of a faithful maniplex is faithful, without
a finiteness assumption on its number of automorphism orbits. For
a faithful thin orientable maniplex, this cover has all three properties.
\end{corollary}
\begin{proof}
Let the universal string group $W$ act on $X$, and let
$N=\bigcap_{x\in X}\Stab_W(x)$ be its action kernel. The orbit of the
tuple $(x)_{x\in X}$ in the product of all rooted copies of $X$ has
stabilizer $N$. It is therefore the regular $W$-set $W/N$.
Projection to a chosen coordinate covers $X$. A regular cover has
a normal stabilizer contained in the stabilizer of that coordinate,
and hence in its core $N$; it thus covers $W/N$. This proves the
minimality assertion. Theorem~\ref{mix} proves the remaining claims.
\end{proof}
The two-factor and finite-orbit cases are results of the source paper.
The arbitrary-family statement records the scope of the same argument;
it is not counted as a newly solved open problem.

\section{Reconstructing a thin poset}
Let $P$ be a finite-rank thin flagged poset. It has a unique minimum
and maximum, every maximal chain meets every rank, and every interval
of rank difference two has exactly two intermediate faces. Its flag
graph $\Fl(P)$ has a unique adjacency of each color and commuting
distant colors. For a proper face $F$, let $\Fl(P)_F$ be the induced
graph on flags containing $F$.

\begin{theorem}\label{reconstruction}
Assume $\Fl(P)$ is connected. The natural rank-preserving map
\[
 \eta:\Pos(\Fl(P))\longrightarrow P
\]
is an isomorphism if and only if $\Fl(P)_F$ is connected for every
proper face $F$. Equivalently, for every such $F$, the flag graphs
of both intervals $[\min P,F]$ and $[F,\max P]$ are connected.
\end{theorem}
\begin{proof}
Deleting color $i$ prevents a flag from changing its rank-$i$ face.
The resulting components are precisely the connected components of
the graphs $\Fl(P)_F$ for rank-$i$ faces $F$. Send each component
to that face. Every face extends to a maximal chain, so this map
is surjective at every rank. It preserves order because intersecting
components have incident labels in $P$. It is bijective exactly when
each graph $\Fl(P)_F$ is connected. In that case, if $F<G$, extend
the pair to a maximal chain; its flag lies in both corresponding
components. Thus the inverse also preserves order.

A maximal chain through $F$ consists independently of a maximal
chain below $F$ and one above $F$. Its allowed adjacencies change
one of these chains at a time. Consequently $\Fl(P)_F$ is the
Cartesian product of the two interval flag graphs. A product of
two nonempty graphs is connected exactly when both factors are
connected: project a path for necessity, and concatenate paths in
the two factors for sufficiency. This proves the equivalent test.
\end{proof}

For finite $P$, the condition is also necessary for an unspecified
isomorphism $\Pos(\Fl(P))\cong P$. A disconnected face
graph splits one face into at least two while every other face has
at least one preimage, increasing that rank's finite face count.
For infinite posets this counting argument does not apply. Accordingly
Theorem~\ref{reconstruction} answers the natural-reconstruction version
of the source's Problem~3, and its unrestricted-isomorphism version
for finite $P$; no stronger interpretation is asserted here.

\section{Attribution and proof boundary}
The source is G.~Cunningham, E.~Moch\'an, A.~Montero, P.~Poto\v cnik
and M.~Toledo, \emph{Faithfulness in maniplexes}, Journal of Algebraic
Combinatorics 64 (2026), article 23,
doi:10.1007/s10801-026-01559-y. Its face equivalence, stabilizer and
transporter criteria, and its two-factor mixing results are recovered
above. Stabilizer cosets and free group actions are classical inputs,
proved here in the form used. The earlier project papers supply the
color-path and covering framework; they remain whole counted common
papers until a separately reviewed consolidation replaces them.
The present addition is the common equivalence-relation formulation,
its arbitrary-family consequence, and the explicit reconstruction
criterion with its stated interpretation. None is assigned a priority
claim. Rank-three classification requires the separate application
and its corrected construction. Finite diagnostic code does not
replace that argument. Lean verification, when supplied, covers only
the declarations listed in the formal scope file, not all graph or
classification assertions in this paper.
\end{document}
